% ANCHOR: end_of_section 1.2
\begin{remark}[Physics notation: the dictionary used in every chapter]
The abstract objects of this manuscript translate into standard physics notation as follows; each later chapter adds the translation of its own objects.
\begin{itemize}
\item Spacetime. $V$ is complexified Euclidean $\mathbb R^4$ with indices $m=1,\dots,4$; the Wick map $v^0=-iv^4$ relates it to Minkowski space, and the Euclidean matrices are $\sigma^m_{\alpha\dot\alpha}=(\vec\sigma,i)$, $\bar\sigma^{m\dot\alpha\alpha}=(-\vec\sigma,i)$, so that $\sigma^\mu v_\mu=\sigma^mv_m$ for Srednicki's $\sigma^\mu=(1,\vec\sigma)$.
\item Four-dimensional spinors. $S_+$ is the undotted Weyl spinor, index $\alpha=1,2$, Srednicki's left-handed $(2,1)$ of $SU(2)_L\times SU(2)_R$; $S_-$ is the dotted one, index $\dot\alpha$, the $(1,2)$. This is the assignment for the orientation $dx^1\wedge dx^2\wedge dx^3\wedge dx^4$, for which $\rho(e_1e_2e_3e_4)=+1$ on $S_+\otimes W$ (Proposition \ref{prop:c1-chirality}). Indices move with $\varepsilon^{12}=\varepsilon^{\dot1\dot2}=+1$, $\varepsilon_{12}=\varepsilon_{\dot1\dot2}=-1$, and $\chi\psi=\chi^\alpha\psi_\alpha$, $\bar\chi\bar\psi=\bar\chi_{\dot\alpha}\bar\psi^{\dot\alpha}$.
\item R-symmetry. $W=\mathbb C^4$ is the fundamental $\mathbf 4$ of $SU(4)_R\cong\mathrm{Spin}(6)$, written with an upper index $w^i$; $W^*$ is the $\bar{\mathbf 4}$, lower index $w_i$. The six transverse directions $W_{\mathbb R}^\perp$ are the vector $\mathbf 6$ of $SO(6)_R$ with components $w_I$, $I=1,\dots,6$, equivalently the antisymmetric $\mathbf 6$ of $SU(4)$: $w_{ij}=w_I\Sigma^I_{ij}$ with the 't Hooft blocks $\Sigma^I$, $w^{ij}=\tfrac12\varepsilon^{ijkl}w_{kl}$, which is the complex conjugate of $w_{ij}$ for real $w$, and $(u,w)=\tfrac14u^{ij}w_{ij}$. The six scalars $\iota_wA$ are $A_I$ in the $SO(6)$ form and $A_{ij}=A_I\Sigma^I_{ij}$ in the $SU(4)$ form, with $\iota_wA=\tfrac14w^{ij}A_{ij}$.
\item Ten-dimensional spinors. $V_{10}=V\oplus W_{\mathbb R}^\perp$ carries the index $M=(m,I)$. $\Sigma$ is a Weyl spinor $\mathbf{16}$ of $\mathrm{Spin}(10)$ and $\Sigma^*$ the other one, $\mathbf{16}'$. Writing an element of $\Sigma$ with a lower index, $\zeta_A$, and an element of $\Sigma^*$ with an upper index, $\chi^A$ ($A=1,\dots,16$), the reduction to $SU(2)_L\times SU(2)_R\times SU(4)$ is
\[
\zeta_A=\big(\xi_\alpha{}^i,\ \bar\xi_{\dot\alpha i}\big)\in(\mathbf 2,\mathbf 1,\mathbf 4)\oplus(\mathbf 1,\mathbf 2,\bar{\mathbf 4}),\qquad \chi^A=\big(\eta^\alpha{}_i,\ \bar\eta^{\dot\alpha i}\big)\in(\mathbf 2,\mathbf 1,\bar{\mathbf 4})\oplus(\mathbf 1,\mathbf 2,\mathbf 4):
\]
on $\Sigma$ the spinor indices are lower and the $SU(4)$ index is upper on the $\mathbf 4$, lower on the $\bar{\mathbf 4}$; on $\Sigma^*$ all positions are reversed. The pairing $\Sigma\otimes\Sigma^*\to\mathbb C$ is $(\zeta,\chi)=\xi_\alpha{}^i\eta^\alpha{}_i-\bar\xi_{\dot\alpha i}\bar\eta^{\dot\alpha i}$, the minus sign on the dotted block being part of the convention of [ESW].
\item Gamma matrices. $\rho(v)=v^M\Gamma_M$ is Clifford multiplication by the ten-dimensional gamma matrices in their chiral blocks, $\Gamma_M\colon\Sigma\to\Sigma^*$ and $\Sigma^*\to\Sigma$, with $\Gamma_M\Gamma_N+\Gamma_N\Gamma_M=2\delta_{MN}$; the $\Gamma$-pairing $\Gamma(\zeta_1,\zeta_2)^M=(\rho(e^M)\zeta_1,\zeta_2)$ is the symmetric bilinear $\zeta_1\Gamma^M\zeta_2$. Along spacetime, on the Dirac column $(\psi_\alpha,\bar\chi^{\dot\alpha})$, $\rho(v)=\gamma_W^\mu v_\mu=-i\gamma_S^\mu v_\mu$ with Weinberg's and Srednicki's gamma matrices after the Wick map; in components $(\rho(v)\xi)^\alpha{}_i=i\,v_m\bar\xi_{\dot\alpha i}\bar\sigma^{m\dot\alpha\alpha}$ and $(\rho(v)\xi)^{\dot\alpha i}=-i\,v_m\bar\sigma^{m\dot\alpha\alpha}\xi_\alpha{}^i$. Along $W_{\mathbb R}^\perp$, $\rho(w)$ is the Yukawa matrix, $(\rho(w)\xi)^\alpha{}_j=\varepsilon^{\alpha\beta}w_{ji}\xi_\beta{}^i$, $(\rho(w)\xi)^{\dot\alpha j}=\varepsilon^{\dot\alpha\dot\beta}w^{ji}\bar\xi_{\dot\beta i}$. On a two-form, $\rho(F)=\rho(q(F))$ is $\tfrac12F_{MN}\Gamma^{MN}$ with $\Gamma^{MN}$ the antisymmetrised product, and ${\not d}_A=\Gamma^MD_M$ reduces to the four-dimensional Dirac operator $\bar\sigma^mD_m$ plus the Yukawa couplings $[A_{ij},-]$.
\item Parities and ghost numbers. $\Pi$ is the parity shift: $\Pi\Sigma\otimes\mathfrak g$ means that the gaugino components $\lambda_\alpha{}^i$, $\bar\lambda_{\dot\alpha i}$ are Grassmann-odd fields of ghost number $0$, four left-handed Weyl fermions in the $\mathbf 4$ and their right-handed partners in the $\bar{\mathbf 4}$; in Euclidean signature $\bar\lambda$ is an independent field, which on a Lorentzian real slice is Srednicki's $\lambda^\dagger$. $\mathfrak g[1]$ places the ghost $c$ in degree $-1$ as a section, so its coordinate has ghost number $+1$. The ghost numbers of the coordinates are $c:+1$, $A,\lambda:0$, $A^*,\lambda^*:-1$, $c^*:-2$, and the Grassmann parities are odd for $c,\lambda,A^*$ and even for $A,\lambda^*,c^*$. The gaugino antifield has components $\lambda^{*\alpha}{}_i$, $\bar\lambda^{*\dot\alpha i}$ in $\Sigma^*$.
\item Gauge algebra and coupling. The coupling sits in the pairing of $\mathfrak g$: for $\mathfrak{su}(N)$, $(X,Y)=-\tfrac{4\hbar}{g^2}\operatorname{tr}(XY)$, and with $A_m=-igA^a_mT^a$, $A_I=-igA^a_IT^a$, $\lambda=-\tfrac{ig}{\sqrt2}\lambda^aT^a$, $\operatorname{tr}(T^aT^b)=\tfrac12\delta^{ab}$, $S/\hbar$ is the Wick image of Srednicki's Lagrangian $-\tfrac14F^{a\mu\nu}F^a_{\mu\nu}-\tfrac12D^\mu A_ID_\mu A_I+i\lambda^{\dagger a}\bar\sigma^\mu D_\mu\lambda^a+\cdots$; the unit $\frac{\hbar\operatorname{tr}_{\mathfrak g}(\mathrm{ad}_X\mathrm{ad}_Y)}{16\pi^2(X,Y)}$ equals $-\frac{g^2N}{32\pi^2}$.
\item Superconformal generators. The odd generators of $\mathfrak{psl}(4|4)$ in the blocks of [SW, (17)] are $Q^i_\alpha$, $\bar Q_{\dot\alpha i}$ (parameters $\xi_\alpha{}^i$, $\bar\xi_{\dot\alpha i}$, an element of $\Sigma$) and $S^\alpha{}_i$, $\bar S^{\dot\alpha i}$ (parameters $\eta^\alpha{}_i$, $\bar\eta^{\dot\alpha i}$, an element of $\Sigma^*$); the even part consists of $P_m$, $K_m$, the dilatation $\Delta$, the rotations $\mathfrak{so}(4)=\mathfrak{su}(2)_L\oplus\mathfrak{su}(2)_R$ and the R-symmetry $\mathfrak{so}(6)=\mathfrak{su}(4)$.
\end{itemize}
\end{remark}
